% status: ZERO
% Printed Chapter 39 / stable source ch22.
\chapter{Speculative Decoding: The Draft--Verify Contract}\label{ch:speculative-decoding}

\begin{ChapterIntent}
Ordinary autoregressive decode pays one target-model step for one token. Speculative
decoding proposes several tokens cheaply, verifies them together with the target, and
commits only the prefix that the target accepts. The speed idea is simple; the
correctness protocol is not. This chapter derives exact speculative sampling, proves
its one-step distribution, traces tentative KV and grammar state through rejection,
and turns proposal, verification, commitment and publication into explicit frontiers.
The result is a reusable contract for n-gram, small-model and later tree-based drafters.
\end{ChapterIntent}

\OpeningSection{Eleven tokens checked; only one history survives}

In this edition's M1 measurement, Llama~3.2~3B was asked to repeat a 218-word passage
already present in its prompt. Ordinary generation produced a 245-token answer at
roughly 15--16 tokens/s. Prompt-lookup speculation proposed 224 tokens; all were
accepted. The answer required roughly 21 target forward passes instead of 245 and ran
at about 60--65 tokens/s: 3.7--4.3$\times$ faster across the recorded runs
(\texttt{research/measurements/2026-09-24-m1-part4.md}). The speculative and ordinary
answers matched character-for-character in those runs.

The gain was workload-dependent. Rewriting the passage with one word changed accepted
193 of 208 proposals and measured 2.5--3.1$\times$ faster. Summarization and open-ended
story generation found no useful repeated 12-token run, proposed nothing, and showed
no gain. A length-limit edge case ended one speculative run early because accepted
draft accounting interacted with the server's output limit.

The measurement captures both sides of speculation. Predictable text lets one target
pass validate many future positions. But those future positions are provisional until
the controller decides which contiguous prefix becomes history.

\Foundations{Proposal, evaluation, commitment and publication are different events}{
A drafter may create tokens the user will never see. The target may evaluate KV for
positions that are later rejected. A verifier may accept several tokens before the
transport publishes them. Correct engines therefore track multiple frontiers instead
of one ``sequence length.'' Exact speculative sampling preserves the target
distribution only when acceptance and rejection correction follow the proper rule.
Chapter~\ref{app:sampling} defines autoregressive sampling; Chapters
\ref{ch:paged-kv}, \ref{ch:structured-output} and \ref{ch:request-lifecycle} provide
the ownership/state machinery that speculation must respect.
}

\NapkinSection{Napkin math: useful tokens per expensive pass}

Let the drafter propose $\gamma$ tokens. In the simplified independent-acceptance
model, each draft is accepted with probability $\alpha$ until the first rejection.
A target pass also yields one correction token after rejection or one bonus token when
all drafts are accepted. The expected emitted tokens are
\begin{equation}
E[T]=1+\alpha+\alpha^2+\cdots+\alpha^\gamma
=\frac{1-\alpha^{\gamma+1}}{1-\alpha}.
\label{eq:spec-etokens}
\end{equation}
If each draft token costs $c$ ordinary target steps and verification costs $v$ target
steps, an approximate speedup is
\begin{equation}
S(\gamma)\approx
\frac{1-\alpha^{\gamma+1}}
{(1-\alpha)(\gamma c+v)}.
\label{eq:spec-speed}
\end{equation}
These are planning equations, not a guarantee: real acceptances are correlated,
verification cost depends on batch/context/hardware, and draft/target execution can
overlap.

The M1 record shows why $v$ matters. A multi-token verification pass did not cost one
single-token decode step; around the measured sizes it cost materially more. A
drafter can achieve excellent acceptance and still lose if verification and draft work
consume the saved time.

\section{Five frontiers, not one length}

For one request define:
\begin{description}
\item[Committed frontier $C$:] tokens already accepted as target history.
\item[Draft frontier $D$:] furthest proposal generated by the drafter.
\item[Evaluated frontier $E$:] furthest position for which target forward state/KV was computed.
\item[Published frontier $P$:] tokens irreversibly exposed to the client/application.
\item[Auxiliary frontier $G$:] grammar/tool/recurrent state corresponding to committed history.
\end{description}
During ordinary decode these often move together. During speculation,
\[
C\le E,\qquad C\le D,
\]
and normally $P\le C$. A transport implementation may buffer several newly committed
tokens before publishing them.

The core invariant is:
\[
\boxed{\text{Only a contiguous verified prefix may advance }C.}
\]
Rejected suffix work may exist beyond $C$ in KV pages, draft-model state, grammar
snapshots or temporary logits. It is garbage until explicitly reclaimed or restored.

\EngineFigure{ch22-frontiers}{One speculative iteration has distinct committed, drafted, target-evaluated and published frontiers. Verification may run ahead; rejection collapses model-derived state back to the accepted boundary before the correction becomes the next history.}{fig:spec-frontiers}

\section{The target defines truth; the drafter defines opportunity}

At a position, let $q$ be the proposal distribution and $p$ the target distribution.
The drafter is an optimization component. It may be an n-gram table, a small language
model, a head attached to the target, or a tree generator. The target distribution is
the semantic contract.

A shortcut such as ``accept whenever the draft probability is high'' changes the
output distribution. So does ``on rejection sample directly from $p$'' when the draft
itself was sampled from $q$. Exact speculative sampling needs a correction for the
probability mass already represented by accepted proposals.

\section{Exact one-token speculative sampling}

Draft $x\sim q$. Accept it with probability
\begin{equation}
a(x)=\min\left(1,\frac{p(x)}{q(x)}\right).
\label{eq:spec-accept}
\end{equation}
If rejected, sample replacement $y$ from
\begin{equation}
r(y)=\frac{[p(y)-q(y)]_+}{\sum_z[p(z)-q(z)]_+},
\qquad [u]_+=\max(0,u).
\label{eq:spec-residual}
\end{equation}
Leviathan et al. and Chen et al. derive this exact correction
\cite{leviathan2023speculative,chen2023speculative}.
The replacement law is the normalized positive residual. Sampling from the
unadjusted target at this point would not preserve the target distribution:
the accepted branch has already supplied some of its probability mass.

Why does it work? The probability of emitting $y$ through the accept branch is
\[
q(y)a(y)=\min(p(y),q(y)).
\]
Let
\[
\beta=\sum_z\min(p(z),q(z))
\]
be total acceptance probability. Since both $p$ and $q$ sum to one,
\[
1-\beta=\sum_z[p(z)-q(z)]_+.
\]
The rejection branch emits $y$ with probability
\[
(1-\beta)r(y)=p(y)-\min(p(y),q(y)).
\]
Adding the two branches gives exactly $p(y)$.

This proof is stronger than observing matching samples. The lab enumerates a finite
distribution and verifies conservation for every token.

\section{A hand calculation}

Let
\[
p=(0.50,0.30,0.20),\qquad q=(0.60,0.10,0.30).
\]
The accept-branch masses are
\[
\min(p,q)=(0.50,0.10,0.20),
\]
so $\beta=0.80$. Positive residual mass is
\[
[p-q]_+=(0,0.20,0),
\]
and rejection always emits token 2. Final mass becomes
\[
(0.50,0.10,0.20)+(0,0.20,0)=(0.50,0.30,0.20)=p.
\]
If an implementation instead sampled from unadjusted $p$ after rejection, the result
would be $(0.60,0.16,0.24)$ for this example, not $p$. The bug can look statistically
reasonable while being mathematically wrong.

\section{Multiple drafts: verify in order}

For drafts $x_1,\ldots,x_\gamma$, the target evaluates distributions for all proposal
positions in one pass. Acceptance still proceeds \emph{sequentially}. At position
$j$, the relevant $p_j$ and $q_j$ are conditioned on the history that includes all
previous accepted drafts.

Stop at the first rejection. Emit its residual correction and discard later proposals.
If all $\gamma$ drafts are accepted, sample a bonus token from the target distribution
at the next position. The target forward pass has therefore converted one expensive
weight traversal into between one and $\gamma+1$ committed tokens.

\EngineFigure{ch22-transaction}{Draft, verify and commit form a transaction. The verifier walks proposals in order; rejection commits only the accepted prefix plus correction and retires the tentative suffix. Full acceptance commits all drafts plus a target bonus.}{fig:spec-transaction}

\section{Deterministic proposals need a different explanation}

Prompt lookup often proposes deterministically rather than sampling from an explicit
$q$. In greedy decoding, verification is straightforward: a proposed token survives
when it equals the target's selected token. That preserves the target's greedy
trajectory subject to ordinary numerical determinism.

For stochastic decoding, do not casually substitute a deterministic guess into the
sampled-$q$ proof above. Engines can implement deterministic proposal verification by
sampling/accepting according to their own exact construction, but the probability
contract must be stated explicitly. ``The draft looks right'' is not a proof.

This distinction matters when reading source: an implementation can be correct without
literally materializing a dense $q$ tensor, but its acceptance path must correspond to
a justified target-preserving algorithm.

\section{Why verification can be cheaper than $\gamma$ decode steps}

Decode is often memory-bandwidth limited because every step streams a large fraction
of model weights for a small amount of arithmetic. Verifying several adjacent
positions resembles a short prefill: the same weights support more matrix work per
load, raising arithmetic intensity.

Verification is not free. Attention touches additional KV positions; temporary
activations/logits grow; kernel shapes change; and large concurrent batches may
already use the accelerator efficiently. Speculation monetizes unused compute only
when the verification shape remains cheaper than the target steps it replaces.

This is why throughput speedup, per-user latency speedup and target-forward-pass
reduction must be reported separately.

\section{Prompt and n-gram lookup}

A model-free drafter searches existing tokens for a suffix matching the current
history, then proposes tokens that followed the earlier occurrence. It is extremely
cheap and excels at quotation, editing, templated output and code with repeated spans.

The opening measurement is almost an ideal case: the answer copies prompt material.
For summarization or creative writing, matching context may disappear and the drafter
naturally proposes little or nothing.

At vLLM commit \texttt{155488d853a0bc42df227dbfc74005b3fd488e94}, V1 contains
\texttt{vllm/v1/spec\_decode/ngram\_proposer.py} and explicit rejection-sampling
paths. SGLang commit \texttt{7a719e9a65e7f52b012ab37211b5f174e5d3d38d}
contains an n-gram worker and separate speculative tree machinery. These are
source-pinned examples, not claims about every release/default.

\section{A smaller model as drafter}

A small model can propose ordinary language beyond repeated spans. If its target-step
cost ratio is $c$, then $\gamma c$ in equation~\eqref{eq:spec-speed} can dominate.
The best drafter is not necessarily the most accurate: it maximizes accepted target
work saved per unit of draft time and memory.

A separate drafter also needs weights, KV/state, tokenizer compatibility and scheduler
capacity. On a memory-constrained device, loading a drafter can reduce the target's KV
capacity enough to lower concurrency. Include that opportunity cost.

Tokenizer mismatch is not merely inconvenience. Draft positions must map to target
token positions for verification. Text-mediated retokenization is possible in some
systems, but it adds conversion complexity and can make one draft token correspond to
multiple target tokens or vice versa.

\section{Tentative KV is a transaction log}

Suppose $C=100$ and the drafter proposes positions 101--104. The target evaluates all
four and writes tentative KV through 104. Verification accepts 101 and 102, rejects
103, and produces correction $z$ for position 103.

KV for 104 is invalid. Depending on implementation, KV for the drafted token at 103
may also not represent correction $z$ and must be recomputed/replaced. The safe state
transition is:
\[
E:104\longrightarrow C:102\longrightarrow \text{commit correction }103.
\]
Paged KV makes truncation cheap when tentative blocks/tail slots are privately owned.
Shared immutable prefix blocks before $C$ remain untouched.

A production engine should not publish a speculative KV block into a shared prefix
cache before its tokens are committed. Otherwise a rejected future can become visible
to another request.

\section{Rollback is representation-dependent}

For token-indexed transformer KV, rollback can lower the valid sequence length and
release tail pages. For recurrent or state-space layers, state at position $t+1$ is a
compressed function of all previous tokens; there may be no per-token state to delete.
The engine needs a checkpoint before speculation or a recomputation strategy.

Thus the speculative controller should depend on a generic state transaction:
\[
\texttt{checkpoint}\rightarrow\texttt{evaluate tentative}
\rightarrow\{\texttt{commit}(k),\texttt{restore}\}.
\]
Chapter~\ref{ch:hybrid-models} will revisit the representation cost.

\section{Structured output adds another transactional state}

A grammar state machine must advance only for committed tokens. One implementation can
clone/checkpoint grammar state, consume proposed tokens tentatively during verification,
then retain only the accepted prefix. Another can delay grammar-state mutation until
the verifier commits.

The verifier also has to mask target and/or proposal distributions consistently with
the structured-output contract. A draft token forbidden by the grammar cannot become
committed merely because the unconstrained target assigns it high probability.

This is the same principle as Chapter~\ref{ch:structured-output}: model state and
grammar state form one logical transaction even when stored separately.

\section{Length limits and stop sequences are commit rules}

A draft pass can accept several tokens at once. If only two output tokens remain under
a request's maximum, the engine must not publish five simply because verification
accepted them. Stop strings and EOS can occur inside the accepted run.

Therefore acceptance and publication are not identical. After mathematical
verification, apply request-level stopping semantics in order and commit/publish only
the legal prefix. State beyond the terminal token is retired.

The opening's early-ending edge case is valuable precisely because it reveals a
boundary where pass accounting and user-visible length accounting can diverge.

\section{Cancellation races with verification}

Cancellation can arrive while the target is verifying a speculative block. Device
work may be non-preemptible, but ownership still changes. Once cancellation is
observed as terminal, no later accepted draft may become user-visible.

Tentative KV and drafter state retire when outstanding work completes. The request
lifecycle's rule still holds:
\[
\text{logical terminality}\ne\text{physical work retirement}.
\]
Speculation increases the amount of orphanable work per in-flight iteration, making
the distinction more important.

\section{Continuous batching makes $\gamma$ a scheduling decision}

In a batch, requests can have different draft lengths and acceptance histories. A
verification request consumes multiple token positions from the scheduler's token
budget. Let request $i$ propose $\gamma_i$ tokens; a simple budget is
\[
\sum_i (1+\gamma_i)\le T_{\rm verify}.
\]
Giving one request a long speculative window can delay ordinary decode tokens for
others. The scheduler can cap $\gamma$ under high load even if a single-user benchmark
prefers a longer draft.

Speculation therefore interacts with fairness, chunked prefill and SLO admission.
Chapter~\ref{ch:when-speculation-loses} develops that load-dependent decision.

\section{Adaptive draft length}

The marginal value of draft $j$ is paid only when the first $j-1$ proposals survive.
Under the simplified model, survival is approximately $\alpha^{j-1}$. Long windows
make sense at high acceptance and cheap verification; they waste work after frequent
early rejection.

A controller can estimate recent acceptance, verification cost and draft cost, then
choose $\gamma$ to maximize predicted useful tokens/time. Avoid adapting on acceptance
alone: a large $\gamma$ can raise acceptance count while lowering wall-clock goodput.

\section{Tree proposals generalize the same contract}

A linear drafter proposes one future. Tree speculation proposes several alternatives
at some depths so target verification has a better chance of containing the eventual
accepted path. The target still defines valid committed history; only one root-to-leaf
path can become the autoregressive sequence.

Tree packing needs position/attention masks that prevent one branch from attending to
siblings. KV ownership becomes branch-aware until verification collapses the tree.
Those mechanisms lead directly to EAGLE and related methods in
Chapter~\ref{ch:modern-speculation}; the transaction contract does not change.

\section{Numerical exactness versus algorithmic exactness}

The rejection proof says the \emph{algorithmic distribution} equals the target
distribution under the same $p$ and sampling semantics. It does not promise that a
different kernel shape computes bit-identical logits.

Verification may batch several positions through GEMMs/attention kernels that use
different reduction orders from one-token decode. Floating-point differences can move
ties or near-ties. Greedy outputs can therefore diverge in rare cases even when the
speculative algorithm is mathematically exact.

A rigorous test suite separates:
\begin{enumerate}
\item probability-law correctness of acceptance/residual sampling;
\item target-logit numerical parity within tolerance;
\item greedy sequence parity under a specified deterministic execution mode;
\item end-to-end service semantics such as limits, cancellation and grammar validity.
\end{enumerate}

\section{Failure injection}

Deliberately implement these wrong variants:
\begin{enumerate}
\item on rejection sample from $p$ instead of normalized $[p-q]_+$;
\item keep KV through the entire evaluated frontier after a rejection;
\item advance grammar state for every draft before verification and never restore it;
\item publish drafts before target acceptance;
\item reuse a speculative tail as a prefix-cache entry before commitment.
\end{enumerate}
Each can produce plausible output. Each violates a distinct invariant.

\section{Observability}

Record proposed, accepted, rejected and bonus tokens; acceptance length distribution;
draft and verify time; verification tokens; target passes avoided; tentative KV bytes;
rollback bytes/blocks; draft method; chosen $\gamma$; and cancellation/terminal
retirement.

Useful ratios include
\[
A_{\rm token}=\frac{N_{\rm accepted}}{N_{\rm proposed}},
\qquad
Y_{\rm pass}=\frac{N_{\rm committed}}{N_{\rm target\ passes}},
\]
but neither is a speed metric. Always correlate them with TTFT, inter-token latency,
throughput and goodput.

\EngineFigure{ch22-state-machine}{Speculative request state distinguishes drafting, target evaluation, ordered verification, commit and rollback. Terminal request ownership can interrupt publication without pretending already-launched device work vanished.}{fig:spec-state-machine}

\EngineFigure{ch22-verify-cost}{The time of one forward pass over $k$ tokens
of one sequence, relative to a single token, on the M1 (measured). A
verification pass over $\gamma$ drafts processes $k=\gamma+1$ tokens; the
dashed line is the assumption behind the textbook speedup formula. On this
machine a multi-token pass runs different matrix kernels from a single-token
one (\Chref{roofline}) and costs 1.4 to 3 single-token steps, with a staircase
at small $k$; beyond about 12 tokens the curve is flat, so a 32-token pass costs
what a 12-token pass does.}{fig:ch22-verify-cost}
\EngineFigure{sys-speculative-sequence}{UML sequence view of one speculative iteration. The controller tests drafts in order and stops at the first rejection. A rejection uses the normalized positive residual of target minus proposal; all accepted drafts permit a bonus sample from the target. Rollback removes tentative rejected state, not the committed prefix.}{fig:sys-speculative-sequence}

\LedgerSection
\begin{ConstraintLedger}
Memory bandwidth & buys & Multi-position verification amortizes target weight traffic when ordinary decode is bandwidth-bound.\\
Compute & spends & Target evaluates rejected proposals and the drafter performs extra work; long windows can waste arithmetic.\\
Memory capacity & spends & Drafter weights/state and tentative target KV consume capacity until commit or rollback.\\
Latency & trades & High acceptance can emit several tokens per target pass; poor acceptance or costly verification increases ITL.\\
Correctness & risks & Wrong residual sampling changes the distribution; stale tentative KV/grammar state corrupts future history.\\
Cost & trades & Fewer target passes can lower serving cost, but reduced batch capacity or extra drafter GPUs can reverse the gain.\\
\end{ConstraintLedger}

\LabSection{prove the sampler; then break the transaction}

The Rust lab in \texttt{code/labs/ch22-speculative/} has two independent layers.
First, it computes the exact finite probability mass of the accept branch and residual
branch for a categorical $p,q$ pair and proves that their sum equals $p$ token by
token. A deliberately wrong ``resample from $p$'' function fails that fixture.

Second, a small speculative-state object tracks committed, drafted, evaluated and
published frontiers. Verification commits only a contiguous accepted prefix and
truncates evaluated state after rejection. Tests reject keeping a stale tail and
publishing beyond commitment. The lab is deterministic and does not claim hardware
speed.

\HappenedSection{speculation became a serving subsystem}

The original speculative decoding/sampling papers established target-preserving
draft-and-verify algorithms \cite{leviathan2023speculative,chen2023speculative}.
Production engines subsequently added multiple proposer families and specialized
verification/sampling paths.

For a source-pinned snapshot, vLLM commit
\texttt{155488d853a0bc42df227dbfc74005b3fd488e94} contains a V1 n-gram proposer and
rejection samplers under its speculative-decode/worker paths. SGLang commit
\texttt{7a719e9a65e7f52b012ab37211b5f174e5d3d38d} contains n-gram and tree-oriented
speculative machinery. The exact current defaults and supported combinations are
moving implementation facts and should be pinned whenever benchmarked.

The durable architecture is proposer $\rightarrow$ target verification $\rightarrow$
ordered acceptance $\rightarrow$ state commit/rollback.

\FrontierSection{2026-10-05}

Linear draft models and n-gram lookup are now only the foundation. Current systems use
target-derived draft heads, tree proposals, EAGLE-family methods, multi-token prediction
and increasingly parallel drafting. Those mechanisms seek higher accepted tokens per
verification without paying a second full autoregressive model.

The next chapter treats those modern drafters. This chapter deliberately freezes the
contract underneath them: regardless of how proposals are generated, target semantics,
ordered verification, state transactions and service terminality remain authoritative.

\ExercisesSection
\begin{enumerate}
\item \emph{Proof.} For $p=(0.4,0.35,0.25)$ and $q=(0.2,0.5,0.3)$, compute
$\beta$, the residual distribution and final emitted probability of every token.
\item \emph{Derivation.} Evaluate equation~\eqref{eq:spec-speed} for
$\alpha=0.7$, $c=0.04$, $v=1.3$, and $\gamma=1,\ldots,8$. Which $\gamma$ wins under
the simplified model?
\item \emph{State trace.} Start at committed/evaluated frontier 80. Draft five
positions, accept the first three, reject the fourth and emit a correction. State the
legal committed and target-KV frontiers before the next iteration.
\item \emph{Implementation.} Add a grammar frontier to the lab. Acceptance criteria:
rejected drafts leave no grammar mutation; committed drafts advance grammar exactly
once; cancellation prevents publication; rollback remains deterministic.
\item \emph{Concurrency.} Explain why maximizing one request's $\gamma$ can reduce
server goodput at high concurrency even when its personal acceptance rate is high.
\item \emph{Falsification.} Construct $p,q$ for which the incorrect ``sample from
$p$ after rejection'' algorithm differs substantially from $p$, and demonstrate the
difference analytically rather than with random sampling.
\end{enumerate}

\Needspace{18\baselineskip}
\section{Worked problems}

\Needspace{12\baselineskip}
\subsection{A finite draft-budget sweep}
\textbf{Problem.} Evaluate the simplified speedup model for
$\alpha=0.7$, $c=0.04$, $v=1.3$ and draft lengths one through eight.

\textbf{Solution.} Use
$S(\gamma)=(1-\alpha^{\gamma+1})/[(1-\alpha)(c\gamma+v)]$.
The derived speedups are approximately
$(1.269,1.587,1.784,1.899,1.961,1.986,1.988,1.975)$.
Length seven narrowly wins for these illustrative inputs. This is not a
measured optimum: changing verification cost with shape, context or load
can change the ordering.

\Needspace{12\baselineskip}
\subsection{Residual proof fixture}
\textbf{Problem.} \emph{Proof.} For $p=(0.4,0.35,0.25)$ and $q=(0.2,0.5,0.3)$, compute
$\beta$, the residual distribution and final emitted probability of every token.

\textbf{Solution.} For $p=(0.4,0.35,0.25)$ and $q=(0.2,0.5,0.3)$,
$\min(p,q)=(0.2,0.35,0.25)$ and $\beta=0.8$. The positive residual is
$(0.2,0,0)$, so rejection emits token 1 with probability one. Final mass is
$(0.2,0.35,0.25)+0.2(1,0,0)=p$.

\Needspace{12\baselineskip}
\subsection{Frontier trace}
\textbf{Problem.} \emph{State trace.} Start at committed/evaluated frontier 80. Draft five
positions, accept the first three, reject the fourth and emit a correction. State the
legal committed and target-KV frontiers before the next iteration.

\textbf{Solution.} From 80, five drafts evaluate through 85. If positions 81--83 are accepted and 84
rejects, the speculative tail cannot survive. Commit 81--83, emit the correction for
84, and ensure the next target state represents the correction rather than the rejected
draft. The legal history is through 84; any evaluated state for rejected 84/85 must be
discarded or overwritten according to the engine's cache convention.
In this lab, committed history ends at 84 but valid target KV ends at 83.
Feeding the correction at 84 extends that KV frontier; publishing its token
does not itself perform that forward computation.

\Needspace{12\baselineskip}
\subsection{Concurrency}
\textbf{Problem.} \emph{Concurrency.} Explain why maximizing one request's $\gamma$ can reduce
server goodput at high concurrency even when its personal acceptance rate is high.

\textbf{Solution.} A long speculative window consumes multiple verification-token budget units and more
temporary KV. At high concurrency it can delay many one-token decode requests. Personal
tokens/pass can rise while queueing and aggregate goodput worsen.

\Needspace{12\baselineskip}
\subsection{Wrong rejection sampler}
\textbf{Problem.} \emph{Falsification.} Construct $p,q$ for which the incorrect ``sample from
$p$ after rejection'' algorithm differs substantially from $p$, and demonstrate the
difference analytically rather than with random sampling.

\textbf{Solution.} For $p=(0.5,0.3,0.2)$, $q=(0.6,0.1,0.3)$, acceptance mass is
$(0.5,0.1,0.2)$ and rejection probability 0.2. Resampling from $p$ incorrectly yields
$(0.6,0.16,0.24)$ rather than $p$.

